% !TEX spellcheck = en_US



% ~~~~~~~~~~~~~~~~~~~~~~~~~~~~~~~~~~~~~~~ V
% (NC) 2026 Haluk Bingol
% github.com/halukbingol/LaTeX-Templates
%
% Licensees may copy, distribute, display, and perform the work and make derivative works 
% and remixes based on it only for non-commercial purposes. 
% ~~~~~~~~~~~~~~~~~~~~~~~~~~~~~~~~~~~~~~~ A




% =====================================================================
%  Git -- Junior Level: Internals, Rewriting History and Automation
%  ---------------------------------------------------------------------
%  Built on hblib/Template/hblibTemplate.tex with the libraries in ../../hblib/.
%  Programs and their real outputs are in codeoutput/:
%     codeoutput/<Name>.java  a program (classes without main() are only compiled)
%     codeoutput/<Name>.sh    a terminal session, run in a copy of codeoutput/<Name>/
%     codeoutput/<Name>.txt   what it printed, made by:  bash run_examples.sh
% =====================================================================

\documentclass[aspectratio=169,10pt,compress]{beamer}

% ~~~~~~~~~~~~~~~~~~~~~~~~~~~~~~~~~~~~~~~~~~~~~~~~~~~~~~~~~~~~~~~~~ V hblib
	\usepackage{../../hblib/hblibPresentation} % lib presentation: theme, i/N, \hSection, algorithm2e, ...
	\usepackage{../../hblib/hblibCodeoutput} % lib code + output: \lstcode, \lstterminal, ...
	\usepackage{../../hblib/hblibDatastructures} % lib data structures: \hString, \hStack, \hQueue
% ~~~~~~~~~~~~~~~~~~~~~~~~~~~~~~~~~~~~~~~~~~~~~~~~~~~~~~~~~~~~~~~~~ A hblib


% xcolor, array (and the L{width} column), listings, tikz and algorithm2e come from the hblib libraries
\usepackage[T1]{fontenc}
\usepackage{lmodern}
\usepackage{booktabs}

\newcommand{\win}{\textbf{Windows}}
\newcommand{\mac}{\textbf{macOS}}

% algorithm2e set-up (the package is loaded by hblibPresentation)
\SetAlCapFnt{\small}
\SetAlCapNameFnt{\small}
\SetKwProg{Fn}{function}{}{end}
\SetKw{KwError}{error}

% ---------------------------------------------------------------------
\title{Git}
\subtitle{Junior Level: Internals, Rewriting History and Automation}
\author[Bingol]{Haluk O. Bingol}

\institute[FCIS]{
    Faculty of Computer and Information Sciences\\
    Yeditepe University
}
\date{
	{\tiny \hbTimeStamp}
}

\begin{document}




% ~~~~~~~~~~~~~~~~~~~~~~~~~~~~~~~~~~~~~~ V frm
\begin{frame}[t,fragile]
  \titlepage
\end{frame}
% ~~~~~~~~~~~~~~~~~~~~~~~~~~~~~~~~~~~~~~ A frm




% ~~~~~~~~~~~~~~~~~~~~~~~~~~~~~~~~~~~~~~ V frm
\begin{frame}[t,fragile] \frametitle{Outline}
  \begin{columns}[T]
    \begin{column}{0.48\textwidth}
      \tableofcontents[sections={1-3}, hideallsubsections]
    \end{column}
    \begin{column}{0.48\textwidth}
      \tableofcontents[sections={4-6}, hideallsubsections]
    \end{column}
  \end{columns}
\end{frame}
% ~~~~~~~~~~~~~~~~~~~~~~~~~~~~~~~~~~~~~~ A frm

% ~~~~~~~~~~~~~~~~~~~~~~~~~~~~~~~~~~~~~~ V
\hSection[Objects]{The Object Database}
% ~~~~~~~~~~~~~~~~~~~~~~~~~~~~~~~~~~~~~~ A




% ~~~~~~~~~~~~~~~~~~~~~~~~~~~~~~~~~~~~~~ V frm
\begin{frame}[t,fragile] \frametitle{Looking inside \texttt{.git/objects}}
  \begin{columns}[T]
    \begin{column}{0.49\textwidth}
      \lstterminal[basicstyle=\ttfamily\tiny, lastline=19]{codeoutput/Objects.txt}
    \end{column}
    \begin{column}{0.49\textwidth}
      \lstterminal[basicstyle=\ttfamily\tiny, firstline=20]{codeoutput/Objects.txt}
    \end{column}
  \end{columns}

  \vspace{0.1em}
  \small
  \texttt{git add} already writes the \textbf{blob}; its name is the SHA-1 of \texttt{"blob 8\textbackslash 0Merhaba\textbackslash n"}, so
  \texttt{hash-object} gives the same name. The copy adds \textbf{no} new object: same content, same blob.
\end{frame}
% ~~~~~~~~~~~~~~~~~~~~~~~~~~~~~~~~~~~~~~ A frm




% ~~~~~~~~~~~~~~~~~~~~~~~~~~~~~~~~~~~~~~ V frm
\begin{frame}[t,fragile] \frametitle{From a commit to the bytes of a file}
  \begin{columns}[T]
    \begin{column}{0.56\textwidth}
      \begin{center}
      \begin{tikzpicture}[ob/.style={draw, thick, rounded corners=2pt, minimum width=2.2cm, minimum height=0.7cm, font=\scriptsize, align=center}]
        \node[ob, fill=gray!20]  (h) at (0,1.4)   {\texttt{HEAD}\\ \texttt{ref: refs/heads/main}};
        \node[ob, fill=colorHead!50] (m) at (3.2,1.4) {\texttt{refs/heads/main}\\ commit hash};
        \node[ob, fill=red!15]   (c) at (3.2,0)   {commit\\ tree, parent, author};
        \node[ob, fill=blue!12]  (t) at (3.2,-1.4) {tree\\ \texttt{hello.txt} $\to$ blob};
        \node[ob, fill=green!15] (b) at (0,-1.4)  {blob \texttt{41979c4}\\ \texttt{Merhaba}};
        \draw[-latex, thick] (h) -- (m);
        \draw[-latex, thick] (m) -- (c);
        \draw[-latex, thick] (c) -- (t);
        \draw[-latex, thick] (t) -- (b);
      \end{tikzpicture}
      \end{center}
    \end{column}
    \begin{column}{0.42\textwidth}
      \begin{itemize}\small
        \item \textbf{Refs} (branches, tags, \texttt{HEAD}) are tiny text files with a hash or another ref's name.
        \item \textbf{Objects} are compressed files named by their hash: immutable.
        \item Every Git command is a walk over this graph; \texttt{cat-file} and \texttt{rev-parse} let you do it by hand.
      \end{itemize}
    \end{column}
  \end{columns}
\end{frame}
% ~~~~~~~~~~~~~~~~~~~~~~~~~~~~~~~~~~~~~~ A frm




% ~~~~~~~~~~~~~~~~~~~~~~~~~~~~~~~~~~~~~~ V frm
\begin{frame}[t,fragile] \frametitle{Refs and \texttt{HEAD}}
  \begin{columns}[T]
    \begin{column}{0.49\textwidth}
      \lstterminal[basicstyle=\ttfamily\tiny, lastline=12]{codeoutput/Refs.txt}
    \end{column}
    \begin{column}{0.49\textwidth}
      \lstterminal[basicstyle=\ttfamily\tiny, firstline=13]{codeoutput/Refs.txt}
    \end{column}
  \end{columns}

  \vspace{0.1em}
  \small
  \texttt{HEAD~1}: the parent; \texttt{HEAD\textasciicircum\{tree\}}: its snapshot. On a branch, \texttt{HEAD} holds \texttt{ref: refs/heads/\dots};
  detached, it holds a commit hash.
\end{frame}
% ~~~~~~~~~~~~~~~~~~~~~~~~~~~~~~~~~~~~~~ A frm

% ~~~~~~~~~~~~~~~~~~~~~~~~~~~~~~~~~~~~~~ V
\hSection[Reset]{Moving Branches: reset and reflog}
% ~~~~~~~~~~~~~~~~~~~~~~~~~~~~~~~~~~~~~~ A




% ~~~~~~~~~~~~~~~~~~~~~~~~~~~~~~~~~~~~~~ V frm
\begin{frame}[t,fragile] \frametitle{Three kinds of reset}
  \begin{center}\small
  \begin{tabular}{llll}
    \toprule
    \texttt{git reset \dots{} C} & moves the branch to \texttt{C} & staging area & working tree \\
    \midrule
    \texttt{-{}-soft}            & yes                            & unchanged     & unchanged \\
    \texttt{-{}-mixed} (default) & yes                            & reset to \texttt{C} & unchanged \\
    \texttt{-{}-hard}            & yes                            & reset to \texttt{C} & \textbf{reset to \texttt{C}}: edits lost \\
    \bottomrule
  \end{tabular}
  \end{center}
  \vspace{0.3em}
  \begin{itemize}\small
    \item \texttt{reset -{}-soft HEAD\textasciitilde1}: ``uncommit'', keep everything staged; useful to redo the last commit.
    \item \texttt{reset -{}-mixed}: unstage. \texttt{reset -{}-hard}: throw away uncommitted work. \hRemark{Uncommitted changes cannot be recovered.}
    \item \texttt{reset} rewrites the branch: do not use it on commits already pushed; use \texttt{revert} there.
  \end{itemize}
\end{frame}
% ~~~~~~~~~~~~~~~~~~~~~~~~~~~~~~~~~~~~~~ A frm




% ~~~~~~~~~~~~~~~~~~~~~~~~~~~~~~~~~~~~~~ V frm
\begin{frame}[t,fragile] \frametitle{reset in action}
  \begin{columns}[T]
    \begin{column}{0.49\textwidth}
      \lstterminal[basicstyle=\ttfamily\tiny, lastline=11]{codeoutput/Reset.txt}
    \end{column}
    \begin{column}{0.49\textwidth}
      \lstterminal[basicstyle=\ttfamily\tiny, firstline=12]{codeoutput/Reset.txt}
    \end{column}
  \end{columns}

  \vspace{0.1em}
  \small
  \texttt{-{}-soft}: the change of \texttt{D} is staged again (\texttt{M\ }); \texttt{-{}-mixed}: unstaged (\texttt{\ M}); \texttt{-{}-hard}: gone.
\end{frame}
% ~~~~~~~~~~~~~~~~~~~~~~~~~~~~~~~~~~~~~~ A frm




% ~~~~~~~~~~~~~~~~~~~~~~~~~~~~~~~~~~~~~~ V frm
\begin{frame}[t,fragile] \frametitle{Lost commits are not lost: \texttt{reflog}}
  \begin{columns}[T]
    \begin{column}{0.49\textwidth}
      \lstterminal[basicstyle=\ttfamily\tiny, lastline=13]{codeoutput/Reflog.txt}
    \end{column}
    \begin{column}{0.49\textwidth}
      \lstterminal[basicstyle=\ttfamily\tiny, firstline=14]{codeoutput/Reflog.txt}
    \end{column}
  \end{columns}

  \vspace{0.1em}
  \small
  The \textbf{reflog} records every position of \texttt{HEAD} (entries kept 90 days by default, 30 for unreachable commits). Commits that no branch points to stay in the database until
  garbage collection; \texttt{HEAD@\{1\}} is ``where HEAD was one step ago''.
\end{frame}
% ~~~~~~~~~~~~~~~~~~~~~~~~~~~~~~~~~~~~~~ A frm




% ~~~~~~~~~~~~~~~~~~~~~~~~~~~~~~~~~~~~~~ V frm
\begin{frame}[t,fragile] \frametitle{Commits on a detached HEAD}
  \begin{columns}[T]
    \begin{column}{0.49\textwidth}
      \lstterminal[basicstyle=\ttfamily\tiny, lastline=11]{codeoutput/Detached.txt}
    \end{column}
    \begin{column}{0.49\textwidth}
      \lstterminal[basicstyle=\ttfamily\tiny, firstline=12]{codeoutput/Detached.txt}
    \end{column}
  \end{columns}

  \vspace{0.1em}
  \small
  Git warns when you leave commits behind and prints their hash; a branch on that commit saves them.
\end{frame}
% ~~~~~~~~~~~~~~~~~~~~~~~~~~~~~~~~~~~~~~ A frm

% ~~~~~~~~~~~~~~~~~~~~~~~~~~~~~~~~~~~~~~ V
\hSection[Rewrite]{Rewriting Local History}
% ~~~~~~~~~~~~~~~~~~~~~~~~~~~~~~~~~~~~~~ A




% ~~~~~~~~~~~~~~~~~~~~~~~~~~~~~~~~~~~~~~ V frm
\begin{frame}[t,fragile] \frametitle{\texttt{fixup} commits and \texttt{rebase -{}-autosquash}}
  \begin{columns}[T]
    \begin{column}{0.49\textwidth}
      \lstterminal[basicstyle=\ttfamily\tiny, lastline=18]{codeoutput/Autosquash.txt}
    \end{column}
    \begin{column}{0.49\textwidth}
      \lstterminal[basicstyle=\ttfamily\tiny, firstline=19]{codeoutput/Autosquash.txt}
    \end{column}
  \end{columns}

  \vspace{0.1em}
  \small
  \texttt{-{}-fixup=C} makes a commit ``\texttt{fixup! C}''; \texttt{rebase -i -{}-autosquash} moves it next to \texttt{C} and melts it in.
  Here \texttt{GIT\_SEQUENCE\_EDITOR} only prints the todo list instead of opening an editor.
\end{frame}
% ~~~~~~~~~~~~~~~~~~~~~~~~~~~~~~~~~~~~~~ A frm




% ~~~~~~~~~~~~~~~~~~~~~~~~~~~~~~~~~~~~~~ V frm
\begin{frame}[t,fragile] \frametitle{Interactive rebase: the todo list}
  \begin{columns}[T]
    \begin{column}{0.48\textwidth}


% +++++++++++++++++++++++++++++++++++++++ V lst
%: -Lst.
\begin{lstlisting}[style=codesnippet, language=bash, basicstyle=\ttfamily\scriptsize]
pick   1a2b3c4 Add parser
reword 5d6e7f8 Add tests     # edit the message
squash 9a8b7c6 Fix tests     # melt into previous
drop   3c4d5e6 Debug prints  # delete
edit   7f8e9d0 Add docs      # stop to amend
\end{lstlisting}
% +++++++++++++++++++++++++++++++++++++++ A lst


    \end{column}
    \begin{column}{0.48\textwidth}
      \begin{itemize}\small
        \item \texttt{git rebase -i HEAD\textasciitilde5} opens this list in your editor (\win: Notepad++ or VS Code if chosen in the installer; \mac: \texttt{vim} unless \texttt{core.editor} is set).
        \item Reorder lines to reorder commits; change the first word to change what happens.
        \item Every rewritten commit gets a \textbf{new hash}.
        \item \texttt{git rebase -{}-abort} returns to the start.
      \end{itemize}
    \end{column}
  \end{columns}
\end{frame}
% ~~~~~~~~~~~~~~~~~~~~~~~~~~~~~~~~~~~~~~ A frm




% ~~~~~~~~~~~~~~~~~~~~~~~~~~~~~~~~~~~~~~ V frm
\begin{frame}[t,fragile] \frametitle{Copying one commit: \texttt{cherry-pick}}
  \begin{columns}[T]
    \begin{column}{0.49\textwidth}
      \lstterminal[basicstyle=\ttfamily\tiny, lastline=10]{codeoutput/CherryPick.txt}
    \end{column}
    \begin{column}{0.49\textwidth}
      \lstterminal[basicstyle=\ttfamily\tiny, firstline=11]{codeoutput/CherryPick.txt}
    \end{column}
  \end{columns}

  \vspace{0.1em}
  \small
  A bug fixed on \texttt{main} is also needed on the \texttt{release} branch: \texttt{cherry-pick} applies the same change as a new commit.
\end{frame}
% ~~~~~~~~~~~~~~~~~~~~~~~~~~~~~~~~~~~~~~ A frm

% ~~~~~~~~~~~~~~~~~~~~~~~~~~~~~~~~~~~~~~ V
\hSection[Search]{Searching History}
% ~~~~~~~~~~~~~~~~~~~~~~~~~~~~~~~~~~~~~~ A




% ~~~~~~~~~~~~~~~~~~~~~~~~~~~~~~~~~~~~~~ V frm
\begin{frame}[t,fragile] \frametitle{Which commit broke it? \texttt{git bisect}}
  \begin{columns}[T]
    \begin{column}{0.49\textwidth}
      \lstterminal[basicstyle=\ttfamily\tiny, lastline=14]{codeoutput/Bisect.txt}
    \end{column}
    \begin{column}{0.49\textwidth}
      \lstterminal[basicstyle=\ttfamily\tiny, firstline=15]{codeoutput/Bisect.txt}
    \end{column}
  \end{columns}

  \vspace{0.1em}
  \small
  Binary search over history: 7 candidate commits need only 3 tests ($\lceil \log_{2} 7 \rceil$). \texttt{bisect run} calls the script:
  exit code 0 = good, 1--127 = bad (125 = skip).
\end{frame}
% ~~~~~~~~~~~~~~~~~~~~~~~~~~~~~~~~~~~~~~ A frm




% ~~~~~~~~~~~~~~~~~~~~~~~~~~~~~~~~~~~~~~ V frm
\begin{frame}[t,fragile] \frametitle{Searching contents: \texttt{-S}, \texttt{-G}, \texttt{-L}}
  \begin{columns}[T]
    \begin{column}{0.49\textwidth}
      \lstterminal[basicstyle=\ttfamily\scriptsize, lastline=5]{codeoutput/Pickaxe.txt}
    \end{column}
    \begin{column}{0.49\textwidth}
      \lstterminal[basicstyle=\ttfamily\scriptsize, firstline=6]{codeoutput/Pickaxe.txt}
    \end{column}
  \end{columns}

  \vspace{0.1em}
  \small
  \texttt{-S}: commits that change the \textbf{number} of occurrences of a string; \texttt{-G}: whose diff matches a regex; \texttt{-L}: history of a line range.
\end{frame}
% ~~~~~~~~~~~~~~~~~~~~~~~~~~~~~~~~~~~~~~ A frm

% ~~~~~~~~~~~~~~~~~~~~~~~~~~~~~~~~~~~~~~ V
\hSection[Hooks]{Automation with Hooks}
% ~~~~~~~~~~~~~~~~~~~~~~~~~~~~~~~~~~~~~~ A




% ~~~~~~~~~~~~~~~~~~~~~~~~~~~~~~~~~~~~~~ V frm
\begin{frame}[t,fragile] \frametitle{A pre-commit hook}
  \begin{columns}[T]
    \begin{column}{0.49\textwidth}
      \lstterminal[basicstyle=\ttfamily\tiny, lastline=8]{codeoutput/Hook.txt}
    \end{column}
    \begin{column}{0.49\textwidth}
      \lstterminal[basicstyle=\ttfamily\tiny, firstline=9]{codeoutput/Hook.txt}
    \end{column}
  \end{columns}

  \vspace{0.1em}
  \small
  Hooks are scripts in \texttt{.git/hooks}; a non-zero exit stops the action. They are \textbf{not} cloned: share them with
  \texttt{git config core.hooksPath tools/hooks} or a tool such as \texttt{pre-commit}. On \win, Git Bash runs them with its own \texttt{sh}.
\end{frame}
% ~~~~~~~~~~~~~~~~~~~~~~~~~~~~~~~~~~~~~~ A frm




% ~~~~~~~~~~~~~~~~~~~~~~~~~~~~~~~~~~~~~~ V frm
\begin{frame}[t,fragile] \frametitle{Useful hooks}
  \begin{center}\small
  \begin{tabular}{lL{0.6\textwidth}}
    \toprule
    \textbf{hook}         & \textbf{typical use} \\
    \midrule
    \texttt{pre-commit}   & format code, run fast checks (no \texttt{TODO}, no secrets, no large files) \\
    \texttt{commit-msg}   & enforce a message style, e.g.\ ``\texttt{JIRA-123: \dots}'' \\
    \texttt{pre-push}     & run the unit tests before pushing \\
    \texttt{post-merge}   & reinstall dependencies after a pull \\
    server: \texttt{pre-receive} & reject pushes that break rules (GitHub: \emph{branch protection} instead) \\
    \bottomrule
  \end{tabular}
  \end{center}
  \small \texttt{git commit -{}-no-verify} skips client hooks: they help, but they are not security. Enforce rules on the server and in CI.
\end{frame}
% ~~~~~~~~~~~~~~~~~~~~~~~~~~~~~~~~~~~~~~ A frm

% ~~~~~~~~~~~~~~~~~~~~~~~~~~~~~~~~~~~~~~ V
\hSection[Summary]{Summary}
% ~~~~~~~~~~~~~~~~~~~~~~~~~~~~~~~~~~~~~~ A




% ~~~~~~~~~~~~~~~~~~~~~~~~~~~~~~~~~~~~~~ V frm
\begin{frame}[t,fragile] \frametitle{Big ideas}
  \begin{itemize}
    \item Git is a database of immutable, content-addressed objects (blob, tree, commit, tag) plus movable refs.
    \item \texttt{reset} moves a branch; \texttt{-{}-soft}/\texttt{-{}-mixed}/\texttt{-{}-hard} decide what happens to the index and the files.
    \item The reflog remembers where \texttt{HEAD} has been: committed work is very hard to lose.
    \item Rewrite \textbf{local} history freely (amend, fixup, interactive rebase, cherry-pick); never rewrite shared history.
    \item \texttt{bisect run} finds a bad commit in $O(\log n)$ tests; \texttt{-S}, \texttt{-G}, \texttt{-L} search contents.
    \item Hooks automate checks locally; real enforcement belongs on the server and in CI.
  \end{itemize}
\end{frame}
% ~~~~~~~~~~~~~~~~~~~~~~~~~~~~~~~~~~~~~~ A frm




% ~~~~~~~~~~~~~~~~~~~~~~~~~~~~~~~~~~~~~~ V frm
\begin{frame}[t,fragile] \frametitle{Exercises}
  \begin{enumerate}\small
    \item Find the blob of a file of yours with \texttt{git rev-parse HEAD:FILE} and print it with \texttt{cat-file -p}.
    \item Delete a branch with unmerged commits (\texttt{branch -D}) and recover it with the reflog.
    \item Make three commits, then use \texttt{rebase -i} to squash them into one with a good message.
    \item Write a test script for one of your projects and find a bug-introducing commit with \texttt{bisect run}.
    \item Write a \texttt{commit-msg} hook that rejects summary lines longer than 50 characters; share it with \texttt{core.hooksPath}.
  \end{enumerate}
\end{frame}
% ~~~~~~~~~~~~~~~~~~~~~~~~~~~~~~~~~~~~~~ A frm




% ~~~~~~~~~~~~~~~~~~~~~~~~~~~~~~~~~~~~~~ V frm
\begin{frame}[t,fragile]
	\begin{center}
		\vspace{4cm}
		\hRemark{\Huge Questions?}
	\end{center}
\end{frame}
% ~~~~~~~~~~~~~~~~~~~~~~~~~~~~~~~~~~~~~~ A frm



\end{document}
